% comp section 08
\section{Learning the horseshoe posterior, and why we usually do not}\label{comp:sec:horseshoe}\label{gbc:sec:horseshoe}

The Gibbs sampler of Chapter~\ref{ch:horseshoe} fits a horseshoe regression in 0.31 seconds at $n = 200$ and $p = 51$ (Section~\ref{comp:sec:cost}), so a network has to earn its training time\index{generative Bayesian computation!horseshoe posterior|textbf}. In our book \emph{Generative Bayesian Computation} (Polson and Sokolov 2026a) we gave it a deliberately easy case: ten coefficients, a fixed design of $n = 50$ rows of independent standard normals, known error standard deviation $\sigma = 0.5$, and 80,000 draws from the prior with the data they imply. The least squares estimate $b$ is sufficient for the coefficient vector $\beta$ given the design and $\sigma$, so each network maps $b$ to the quantile function of one $\beta_j$. In this experiment $\beta_j \mid \lambda_j, \gscale \sim N(0, \lambda_j^2\gscale^2)$ with standard half Cauchy local and global scales, not multiplied by $\sigma$ as in Chapter~\ref{ch:horseshoe}. Each network has width 256 and was trained for 3,000 full batch epochs with the defaults of the \texttt{gbc} package, whose loss adds an L1 location anchor and a crossing penalty to the check loss; the reference Gibbs sampler kept 20,000 draws after discarding 2,000.

The first obstacle is scale. The horseshoe prior has tails at least as heavy as a Cauchy's, so under it $\beta_j$ has no mean and no variance. The sample standard deviations of the 80,000 draws run from 808 to 21,740, while the median of $|\beta_j|$ is 0.56. Standardizing by the training mean and standard deviation squashes 92 percent of the draws to within 0.01 of the mean on the standardized scale, and an error of 0.001 on the standardized scale is an error of at least 0.81 in $\beta_j$, more than the median of $|\beta_j|$.

Three pipelines differ only in scale. The raw pipeline maps $b$ to $\beta_j$. The asinh pipeline maps $\operatorname{asinh}(b)$ to $\operatorname{asinh}(\beta_j)$ and returns the quantiles of $\beta_j$ through $\sinh$. The residual pipeline maps $\operatorname{asinh}(b)$ to $\beta_j - b_j$ and adds $b_j$ back to the fitted quantiles. Only the residual target is easy to scale: given $\beta$, the error $b - \beta$ is normal with covariance $\sigma^2(X^\top X)^{-1}$, which does not depend on $\beta$, so $\beta_j - b_j$ is exactly normal under the prior predictive distribution, with standard deviations from 0.067 to 0.096 across the coefficients.

We fixed two criteria before the runs. A pipeline is calibrated if the central 90 percent interval of each of its ten networks contains $\beta_j$ in 85 to 95 percent of 2,000 new data sets. It agrees with the Gibbs sampler of Makalic and Schmidt (2016), whose own intervals, pooled over ten coefficients and 500 data sets, covered 89.9 percent, if for every coefficient the median over 500 data sets of the Kolmogorov and Smirnov distance between the two posteriors is below 0.05. No pipeline passed. Two raw networks returned intervals of zero width that never contained $\beta_j$, and the other eight were 17 to 88 times as wide as the sampler's. The asinh and residual pipelines were calibrated and still disagreed with the sampler, by median distances of 0.16 to 0.18 and 0.062 to 0.078. Coverage averaged over data sets can be right while the posterior for a given data set is wrong. The residual intervals came within 5 percent of the sampler's widths, but their posterior medians were only 1 to 7 percent closer to the truth than the least squares estimates.

The sampler was also far cheaper: 40.6 seconds on one core for the 500 data sets, run 100 at a time by a vectorized sampler, or 0.081 seconds each, while one network took more than an hour to train, so the ten networks of one pipeline cost what the sampler spends on more than 440,000 data sets. 

Training a coupled horseshoe network would mean drawing the local and global scales from half Cauchy distributions, shared between the treatment and outcome equations when the coupling is through scales; leaving $\tau$ and the coefficient of the propensity covariate unshrunk, with wide normal priors; and holding the design at the observed covariates, so that one network serves one study. With $p$ larger than $n$ the least squares estimate is not unique, but for a fixed design $X^\top Y$ is still sufficient for the coefficients given $\sigma$ (with $Y^\top Y$ added when $\sigma$ is unknown); it has $p$ coordinates, and reducing it brings back the choices of Section~\ref{comp:sec:summaries}. We have not tried a network on the coupled design of Chapters~\ref{ch:dml} and~\ref{ch:horseshoe} ($n = 200$, $p = 51$, five confounders).
